% Physical pages 3--4 of the concise study: the thematic orientation, in the
% order of the Mass. Every claim and quotation comes from the reviewed
% expansive study; notes.md beside this file gives the source lines.
\runninghead{Themes and Movement}
\section{The Propers: Themes and Movement}

The Mass of Christ the King hails as King a Lamb that was slain and a prisoner
bound before Pilate. In the Fathers and Doctors who expound its texts, his
kingdom is the people he has redeemed: the Son, King by nature, receives them
as man by asking and by dying, gathers them out of every nation by the truth,
and hands them over to the Father without ceasing to reign.

\subsection{Introit, Collect and Epistle: the Lamb who receives}

The Introit comes from the fifth chapter of the Apocalypse. The Lamb has taken
the sealed book that no one else was worthy to open, and the elders give the
reason: ``thou wast slain and hast redeemed us to God, in thy blood.'' Then
the angels sing: ``The Lamb that was slain
is worthy to receive power and divinity and wisdom and strength and honour.''
If the Lamb is the eternal Son, how can he receive divinity?

The Fathers answer from the Son's two natures. St Hilary of Poitiers says what
the risen Christ asked for and was given: \latin{ut esset, quod fuerat ante},
that he might be what he had been before. The servant the Son became is
raised into the glory the Word always had. Cassiodorus makes it a rule:
whenever Christ asks that something be given him, his humanity is shown, for
\latin{natura humanitatis accepit, quod semper divina possedit}, the human
nature received what the divine always possessed.

The Collect confesses that God willed to restore all things in his beloved
Son, \latin{universorum Rege}, the King of all, and asks that the families of
the nations, \latin{peccati vulnere disgregatae}, scattered by the wound of
sin, be brought under his \latin{suavissimo \dots\ imperio}, his most gentle
rule. The Epistle (Col 1:12--20) thanks the Father ``who hath
delivered us from the power of darkness and hath translated us into the
kingdom of the Son of his love''; St Thomas Aquinas glosses the verb
\latin{ut essemus regnum Dei}, that we might be the kingdom of God.

\subsection{Gradual, Alleluia and Gospel: a kingdom not from hence}

The Gradual sings from Psalm~71 that ``he shall rule from sea to sea, and from
the river unto the ends of the earth \dots\ all kings of the earth shall adore
him: all nations shall serve him.'' The Alleluia sings Daniel's vision of one
like a son of man: ``his power is an everlasting power that shall not be taken
away: and his kingdom that shall not be destroyed.''

The Gospel (Jn 18:33--37) then sets this King in the governor's hall, bound
and handed over by his own nation. Under Pilate's questions, Jesus says: ``My kingdom is not of this world. If my kingdom were of this
world, my servants would certainly strive that I should not be delivered to
the Jews: but now my kingdom is not from hence.'' How can a kingdom not from
hence claim the kings of the Gradual?

St Augustine answers from the exact words. Christ did not say that his
kingdom is not in this world, but that it is not of it; nor that it is not
here, but that it is not from hence. ``For His kingdom is here until the end
of the world, having tares intermingled therewith until the harvest,'' and
yet ``it only sojourns as a stranger in the world.'' The kingdom is ``those
who believe in Him,'' and they are gathered by the Gospel's last words,
``Every one that is of the truth heareth my voice,'' where to hear is to
believe and obey. St John Chrysostom takes the kingdom for the royal power. It
is not of this world ``because He has his empire from above,'' and the
servants who do not fight are the proof: an earthly king's ``strength lies in
servants; but that which is above is sufficient for itself, needing
nothing.'' Aquinas keeps both senses and concludes \latin{non est hinc \dots\
est tamen hic, quia ubique est}: it is not from here, and yet it is here,
because it is everywhere.

Defending the imperial laws against the Donatists, Augustine reads the second
psalm's ``Serve ye the Lord with fear'' with the Gradual's verse: ``a man
serves God in one way in that he is man, in another way in that he is also
king,'' as a man by living faithfully, as a king by his laws. St Robert Bellarmine
finds the kings' adoration fulfilled in part as nations are converted, ``to
the letter'' at the judgment, and ``always'' in Christ's power over every
ruler.

\subsection{Offertory and Secret: ask of me}

As the gifts are brought to the altar, the Offertory sings the Father's word
to his Son in the second psalm: ``Ask of me, and I will give thee the Gentiles
for thy inheritance, and the utmost parts of the earth for thy possession.''
In the psalm it follows the decree ``Thou art my son, this day have I begotten
thee,'' and a Son so addressed would seem to lack nothing.

Hilary finds the asking in Christ's prayer at the Supper, ``Father, the hour
is come. Glorify thy Son'': \latin{accepit ergo hereditatem gentium, quam
poposcit}, he received the inheritance of the nations, which he asked for.
Augustine hears the command to ask spoken to the Son made man, ``who offered
up Himself as a Sacrifice in the stead of all sacrifices.'' Aquinas gives
Christ the nations \latin{duplici jure}, by a double right: by inheritance,
through his eternal birth, and by merit, through his prayer and his passion.
God wills that even what is owed be received by asking. Chrysostom, on ``For
this was I born,'' objects that born a king, the Son ``has nothing which He
received in addition'': he refers every text of giving to the eternal
generation, where Cassiodorus refers it to the humanity.

The Secret offers \latin{hostiam \dots\ humanae reconciliationis}, the victim
of mankind's reconciliation, \latin{quem sacrificiis praesentibus immolamus},
the Son himself, immolated in the sacrifice now offered. The Collect prayed
for nations scattered by sin; the Secret asks for all nations \latin{unitatis
et pacis dona}, the gifts of unity and peace. Taken with Augustine's account
of the psalm, the sacrifice by which Christ asks for the nations is the one
offered on this altar.

\subsection{Preface, Communion and Postcommunion: the kingdom handed over}

The Preface gives thanks that the Father anointed his only-begotten Son
eternal Priest and King of all things, so that, offering himself on the altar
of the cross, he might accomplish the mysteries of man's redemption, and, all
creatures being subjected to his rule, deliver to the Father ``an eternal and
universal kingdom, a kingdom of truth and life; a kingdom of holiness and
grace, a kingdom of justice, love and peace.'' Its ``deliver'' is St
Paul's ``delivered up'' (1 Cor 15:24); yet the Alleluia sang a power ``that shall not
be taken away.''

Augustine asks what the kingdom is. To deliver it up is ``as much as to say,
When He shall have brought believers to the contemplation of God, even the
Father.'' Hilary argues from the wording: ``It is not said, He shall deliver
up His Kingdom, but, He shall deliver up the Kingdom, that is, deliver up to
God us who have been made the Kingdom by the glorifying of His body.'' Both
Fathers cite, for the purity this sight requires, ``Blessed are the pure in
heart, for they shall see God.''

The Communion sings from Psalm~28 that ``the Lord shall sit king for ever
\dots\ the Lord will bless his people with peace,'' leaving out the psalm's ``The Lord maketh
the flood to dwell.'' For Augustine the flood is this world, where the Lord dwells first ``in
His Saints, kept safely in the Church, as in the ark,'' and ``afterward He
will sit reigning in them for ever.'' Read with his account of the handing
over, the King goes on reigning in the saints he brings to the Father.

The Postcommunion, after \latin{immortalitatis alimoniam}, the food of
immortality, asks that those who glory \latin{sub Christi Regis vexillis
militare}, to fight under the standards of Christ the King, may reign with him
for ever. The Gospel's King has servants who do not fight for him; Hilary gives
this warfare its sense: the kingdom of God
is where sin is overcome and death destroyed, and the co-heirs of the eternal
King must themselves become kings. Blessed Ildefonso Schuster names the
standard: ``The Royal Standard of Jesus Christ is the Cross: \latin{Dicite in
nationibus: regnavit a ligno Deus},'' tell it among the nations, God has
reigned from the tree, ``for the kingdom of Christ is that of obedience,
humility and sacrifice.''
